\documentclass[10pt,a4paper]{amsart}
\usepackage[margin=19mm]{geometry}
\usepackage{amsmath,amssymb,mathtools}
\usepackage[scr=rsfs,cal=euler]{mathalfa}
\usepackage{xfrac}
\usepackage[unicode,hidelinks,bookmarks=false]{hyperref}
\newcommand{\ie}{i.e.\ }
\newcommand{\cf}{cf.\ }
\newcommand{\resp}{respectively\ }
\newcommand{\Subsection}{\subsection}
\providecommand{\nobreakdash}{\nobreak\hbox{-}}
\setlength{\emergencystretch}{2em}
\multlinegap0pt
\usepackage{bm}
\usepackage{bbm}
\usepackage{dsfont}
\usepackage{xspace}
\usepackage{cancel}
\usepackage{mathtools}
\usepackage[shortlabels]{enumitem}
\usepackage{amsthm}
\makeatletter
\@ifundefined{theorem}{\newtheorem{theorem}{Theorem}}{}
\@ifundefined{corollary}{\newtheorem{corollary}[theorem]{Corollary}}{}
\makeatother
\newcommand{\aaa}{\mathcal{A} }
\newcommand{\bp}[1]{\left\lbrace #1 \right\rbrace}
\newcommand{\closed}[1]{\mathcal{C}^{\mathcal{A}}_{#1}}
\newcommand{\open}[1]{\mathcal{N}^{\mathcal{A}}_{#1}}
\newcommand{\spa}{\sigma^\aaa_P}
\newcommand{\wall}[2]{\mathsf{W}(#1,#2)}
\title[Reflecting compact $T\_1$-spaces into bounded distributive la]{Reflecting compact $T\_1$-spaces into bounded distributive lattices}
\author{Mai Gehrke, Elena Pozzan, Matteo Viale}
\date{Source: arXiv:2411.13482v3 (CC BY 4.0)}
\begin{document}
\maketitle

\noindent\emph{Source and licence.} Reflecting compact $T_1$-spaces into bounded distributive lattices. Version arXiv:2411.13482v3 (CC BY 4.0), distributed under the Creative Commons Attribution 4.0 International licence, as recorded in the arXiv licence metadata for this exact version and shown on the abstract page https://arxiv.org/abs/2411.13482v3. Full rights notice: licenses/arxiv-2411.13482-CC-BY-4.0.txt.\par\smallskip

\noindent\textit{Editorial scope.} Counted unit: the Theorem whose source label is \texttt{thm:wallmancompactification-lattice} in the article (source0.tex lines 913--916, source label \texttt{thm:wallmancompactification-lattice}), delivered together with the article's own complete original proof of it (source0.tex lines 918--961, one \texttt{proof} environment of 44 lines). No mathematics is written, repaired or re-derived: statement and proof are the article's own text line for line. Required context retained uncounted: cor:idealsandfilters (lines 567--570). Every definition, display and label of those lines is retained unchanged. Omitted from this package and not redistributed: 1--566, 571--912, 917--917, 962--2742. Nothing inside a retained excerpt is omitted or altered: every retained body is delivered exactly as the article's lines are written, so no omission receipt of any kind accompanies this package. Citations: the retained excerpts carry no citation command at all, so no bibliography entry of the article is transcribed and no reference is redistributed. Reference adaptation: none was needed and none was made; every reference inside the delivered unit resolves inside it, so no unresolved reference or citation is produced. Printed slips kept literal: no printed word, symbol, hypothesis or number of the counted statement or of its proof was altered. Layout and class: the article's own class is \texttt{article} with packages including geometry, babel, amsthm, enumitem, amscd, todonotes, amssymb, amsmath, tikz-cd, graphicx, hyperref, stackrel, fixme; the wrapper is the lane's pinned \texttt{amsart} editorial wrapper at 10pt on A4, which restates the theorem-like environments the retained text needs and adds the article's own macro definitions that the retained text needs (\texttt{\textbackslash aaa}, \texttt{\textbackslash bp}, \texttt{\textbackslash closed}, \texttt{\textbackslash open}, \texttt{\textbackslash spa}, \texttt{\textbackslash wall}), with extra wrapper packages (bm, bbm, dsfont, xspace, cancel, mathtools, enumitem). The delivered numbering is the unit's own. Assets: no retained span contains a figure environment, a graphics-inclusion command, a PDF-page inclusion, a table or a TikZ picture; no image file is copied into this package, the package declares no asset dependency and no image file is redistributed with it. Licence: version arXiv:2411.13482v3 of this article is distributed under the Creative Commons Attribution 4.0 International licence, as recorded in the arXiv licence metadata for this exact version and shown on the abstract page https://arxiv.org/abs/2411.13482v3. A full rights notice is delivered at licenses/arxiv-2411.13482-CC-BY-4.0.txt. Characters: every retained excerpt is pure ASCII, so the delivered engine, XeLaTeX with its default Latin Modern text font, reports no missing character.
\begin{corollary}
\label{cor:idealsandfilters}
    Let $L$ be a  bounded distributive lattice. Then  $J$ is maximal ideal of $L$ if and only if $J^C$ is a minimal prime filter.
\end{corollary}

\begin{theorem}
\label{thm:wallmancompactification-lattice}
Let $P$ be a subfit lattice and $\aaa$ be a large  family of minimal prime filters on $P$. Then, $(\wall{\spa}{\aaa},\wall{\spa}{\tau^\aaa_P} )$ is homeomorphic to $(\aaa_P, \tau^{\aaa_P}_P)$ where $\aaa_P$ is the family of \emph{all} the minimal prime filters on $P$.
\end{theorem}

\begin{proof}
Let 
\begin{align*}
    \phi: (\wall{\spa}{\aaa},\wall{\spa}{\tau^\aaa_P} ) &\to (\aaa_P, \tau^{\aaa_P}_P)\\
    \mathcal{F}=\bp{\closed{p}}_{p \in _{I_\mathcal{F}}}&\mapsto I_\mathcal{F}^C
\end{align*} 
\begin{description}[font=\normalfont\itshape]
        \item[$\phi$ is well defined and bijective.]
        We claim that $\mathcal{F}=\bp{\closed{p}}_{p \in I_{\mathcal{F}}}$ is a $\spa$-ultrafilter if and only if $I_{\mathcal{F}}$ is a maximal ideal on $P$. 
        
     It is evident that $\open{q} \subseteq \open{p}$ if and only if $\closed{p}\subseteq \closed{q}$, thus $I_{\mathcal{F}}$ is closed downwards with respect to $\leq$ if and only if $\mathcal{F}$ is closed upwards with respect to inclusion. Note that 
\[
\closed{p} \cap \closed{q} = \closed{p \vee q}
\]
holds (since the filters in $\mathcal{A}$ are prime): therefore $\mathcal{F}$ is closed under finite intersections if and only if $I_{\mathcal{F}}$ is closed under finite joins. 

Additionally, $\emptyset \in \mathcal{F}$ if and only if $1_P \in I_{\mathcal{F}}$.

Clearly, $\mathcal{F}$ is maximal in $\spa$ if and only if $I_{\mathcal{F}}$ is maximal in $P$.

 Thus, by Corollary \ref{cor:idealsandfilters}, we can conclude that $\phi$ is well defined and bijective.  
 
        \item[$\phi$ is  open and continuous.]
We show that 
\[
\phi[\wall{\spa}{\open{p}}]
=\mathcal{N}^{\aaa_P}_p.
\] 


Recall that for $\mathcal{F}$ in $\spa$
\[
I_{\mathcal{F}}=\bp{q: \closed{q}\in\mathcal{F}}
\]
is a maximal ideal and $\phi(\mathcal{F})=I_\mathcal{F}^C$ is a minimal prime filter.

Now $\mathcal{F} \in \wall{\spa}{\open{p}}$ if and only if there is
$q \in P$ such that $\closed{q}\subseteq\open{p}$ and $\closed{q} \in \mathcal{F}$, i.e. if and only if there is $q\in I_\mathcal{F}$ and $q\vee p=1_P$, i.e. if and only if $p \in \phi(\mathcal{F})$, i.e. if and only if
$\mathcal{F} \in \mathcal{N}^{\aaa_P}_p$.



\end{description}
\end{proof}

\end{document}
